\chapter{函数极限的严格定义与基本运算}

\section*{A组　定义、量词与基本性质}
\addcontentsline{toc}{section}{A组　定义、量词与基本性质}

\begin{solution}{14.1}
\(\{1/n:n\in\N\}\) 的唯一聚点是 \(0\)。任意 \(\delta>0\) 可取 \(n>1/\delta\)，使 \(0<1/n<\delta\)。若 \(a\ne0\)，在 \(a<0\) 时可取不与正半轴相交的小邻域；在 \(a>0\) 时，序列 \(1/n\) 在 \(a\) 附近只有有限项，因而可缩小邻域避开它们。

\(\Q\cap[0,1]\) 的聚点集为 \([0,1]\)：有理数在每个非退化实区间中稠密，端点用单侧点即可；区间外的点与 \([0,1]\) 距离为正。\(\Z\) 没有聚点，因为对每个 \(m\in\Z\) 可取半径 \(1/2\) 排除其余整数，非整数点也可取小邻域至多含一个整数。
\end{solution}

\begin{solution}{14.2}
若 \(a\) 是孤立点，存在 \(\delta_0>0\) 使
\[
 D\cap U^\circ(a,\delta_0)=\varnothing.
\]
于是对任意 \(L\) 和任意 \(\varepsilon>0\)，选 \(\delta\le\delta_0\) 后，定义中的蕴含式因前件从不成立而为真。这样每个 \(L\) 都会成为“极限”，与唯一性冲突。聚点条件排除了这种空泛成立。
\end{solution}

\begin{solution}{14.3}
设 \(g(x)=f(x)\) 对所有 \(x\in D\setminus\{a\}\) 成立，而 \(f(a),g(a)\) 可以不同。极限定义只检验 \(0<\abs{x-a}<\delta\) 的点，因此在这些点上
\[
 \abs{g(x)-L}=\abs{f(x)-L}.
\]
同一个 \(\delta(\varepsilon)\) 同时适用于两函数，故二者在 \(a\) 处有相同极限。
\end{solution}

\begin{solution}{14.4}
严格形式为
\[
 \forall\varepsilon>0\ \exists\delta>0\ \forall x\in D,\quad
 0<\abs{x-a}<\delta\Rightarrow\abs{f(x)-L}<\varepsilon.
\]
否定时依次交换全称与存在量词，并否定蕴含式，得到
\[
 \exists\varepsilon_0>0\ \forall\delta>0\ \exists x\in D,\quad
 0<\abs{x-a}<\delta,\quad\abs{f(x)-L}\ge\varepsilon_0.
\]
\end{solution}

\begin{solution}{14.5}
给定 \(\varepsilon>0\)，取 \(\delta=\varepsilon/2\)。若 \(0<\abs{x-3}<\delta\)，则
\[
 \abs{(2x-5)-1}=2\abs{x-3}<2\delta=\varepsilon.
\]
故极限为 \(1\)。
\end{solution}

\begin{solution}{14.6}
分解
\[
 \abs{x^3-a^3}=\abs{x-a}\abs{x^2+ax+a^2}.
\]
先要求 \(\abs{x-a}<1\)，则 \(\abs x\le\abs a+1\)，所以
\[
 \abs{x^2+ax+a^2}
 \le(\abs a+1)^2+\abs a(\abs a+1)+a^2=:C_a.
\]
这里 \(C_a>0\)。给定 \(\varepsilon>0\)，取
\[
 \delta=\min\{1,\varepsilon/C_a\}.
\]
于是目标误差小于 \(\varepsilon\)。
\end{solution}

\begin{solution}{14.7}
反三角不等式给出
\[
 \bigl|\abs{f(x)}-\abs L\bigr|\le\abs{f(x)-L}.
\]
对给定 \(\varepsilon>0\)，沿用 \(f(x)\to L\) 对应的半径，右端小于 \(\varepsilon\)，故 \(\abs{f(x)}\to\abs L\)。
\end{solution}

\begin{solution}{14.8}
在极限定义中取 \(\varepsilon=L/2\)。存在 \(\delta>0\) 使删心邻域内
\[
 \abs{f(x)-L}<L/2.
\]
因而 \(f(x)>L-L/2=L/2\)。条件 \(L>0\) 保证所选误差为正并给出正下界。
\end{solution}

\section*{B组　运算、估计与反例}
\addcontentsline{toc}{section}{B组　运算、估计与反例}

\begin{solution}{14.9}
设在 \(a\) 的某删心邻域内 \(f(x)\le h(x)\le g(x)\)，且 \(f(x),g(x)\to L\)。给定 \(\varepsilon>0\)，取共同半径，使
\[
 L-\varepsilon<f(x),\qquad g(x)<L+\varepsilon
\]
同时成立，再与次序合并：
\[
 L-\varepsilon<h(x)<L+\varepsilon.
\]
故 \(h(x)\to L\)。
\end{solution}

\begin{solution}{14.10}
若题设中的点列 \((x_n)\subset D\setminus\{a\}\) 确实趋于 \(a\)，并且 \(f,g\) 的函数极限都存在，则保序结论仍成立。记极限为 \(L,M\)。若 \(L>M\)，取 \(\varepsilon=(L-M)/3\)。由两个函数极限和 \(x_n\to a\)，充分大的 \(n\) 满足
\[
 f(x_n)>L-\varepsilon>M+\varepsilon>g(x_n),
\]
与 \(f(x_n)\le g(x_n)\) 矛盾。

这条较弱假设只足以比较两个已经存在的极限，并不能保证极限存在。例如 \(f(x)=\sin(1/x)\)、\(g(x)=1\) 在 \(x_n=1/(2\pi n)\) 上满足 \(f(x_n)\le g(x_n)\)，但 \(f\) 在 \(0\) 处无极限。若点列不趋于 \(a\)，则连极限次序也无法推出。
\end{solution}

\begin{solution}{14.11}
有
\[
 \abs{x^2+3x-10}
 =\abs{x-2}\abs{x+5}.
\]
若 \(\abs{x-2}<1\)，则 \(\abs{x+5}\le\abs{x-2}+7<8\)。给定 \(\varepsilon>0\)，取
\[
 \delta=\min\{1,\varepsilon/8\},
\]
便有目标误差小于 \(8\delta\le\varepsilon\)。
\end{solution}

\begin{solution}{14.12}
对 \(x\ge0\)，
\[
 \abs{\sqrt x-2}
 =\frac{\abs{x-4}}{\sqrt x+2}
 \le\frac12\abs{x-4}.
\]
给定 \(\varepsilon>0\)，取 \(\delta=2\varepsilon\)，即得 \(\abs{\sqrt x-2}<\varepsilon\)。
\end{solution}

\begin{solution}{14.13}
设在某删心邻域内 \(\abs{g(x)}\le M\)。若 \(M=0\)，乘积恒为零；若 \(M>0\)，给定 \(\varepsilon>0\)，由 \(f(x)\to0\) 取半径使
\[
 \abs{f(x)}<\varepsilon/M.
\]
再与有界性半径取最小值，便有 \(\abs{f(x)g(x)}<\varepsilon\)。
\end{solution}

\begin{solution}{14.14}
均在 \(x\to0\) 时考虑，取 \(f(x)=x\)。令
\[
 g_1(x)=\frac1{\sqrt{\abs x}},\qquad
 g_2(x)=\frac1x,\qquad
 g_3(x)=\frac{\sin(1/x)}x.
\]
三者在每个删心邻域内无界，而
\[
 f g_1=\operatorname{sgn}(x)\sqrt{\abs x}\to0,\quad
 f g_2=1,\quad
 f g_3=\sin(1/x)
\]
分别趋于零、趋于 \(1\)、无极限。最后一项沿 \(x_n=1/(\pi/2+2\pi n)\) 与 \(y_n=1/(3\pi/2+2\pi n)\) 分别取值 \(1,-1\)。
\end{solution}

\begin{solution}{14.15}
设 \(g(x)\to M\ne0\)。取误差 \(\abs M/2\)，可得某删心邻域内
\[
 \abs{g(x)}>\abs M/2,
\]
故倒数有定义。再由
\[
 \abs{\frac1{g(x)}-\frac1M}
 =\frac{\abs{g(x)-M}}{\abs{g(x)}\abs M}
 \le\frac2{\abs M^2}\abs{g(x)-M}
\]
知倒数趋于 \(1/M\)。非零条件既使 \(1/M\) 有定义，又产生分母的正下界。
\end{solution}

\begin{solution}{14.16}
若 \(q(a)\ne0\)，商法则直接给出有限极限，例如
\[
 \frac{x+1}{x+2}\longrightarrow\frac23\quad(x\to1).
\]
若 \(q(a)=0\) 且分子分母有可约公因子，约去后可能有可去型有限极限，例如
\[
 \frac{x^2-1}{x-1}=x+1\to2\quad(x\to1).
\]
若没有足够的公共零点抵消，可能出现无穷行为或双侧极限不存在，例如
\[
 \frac1{x^2}\to+\infty,\qquad \frac1x
\]
在 \(x\to0\) 时左右符号相反。仅凭 \(q(a)=0\) 不能作统一结论，须比较零点阶或直接估计。
\end{solution}

\begin{solution}{14.17}
令 \(b=0\)、\(g(x)\equiv0\)。定义
\[
 f(y)=
 \begin{cases}
 y,&y\ne0,\\
 1,&y=0.
 \end{cases}
\]
则 \(g(x)\to0\)、\(\lim_{y\to0}f(y)=0\)，但 \(f(g(x))\equiv1\)。内层函数持续取中心值，而外层删心极限不控制该点值。
\end{solution}

\begin{solution}{14.18}
设 \(f\) 在 \(b\) 连续。给定 \(\varepsilon>0\)，存在 \(\eta>0\) 使
\[
 y\in E,\quad\abs{y-b}<\eta
 \Longrightarrow\abs{f(y)-f(b)}<\varepsilon.
\]
这里允许 \(y=b\)，因为此时误差为零。由 \(g(x)\to b\) 取 \(\delta\)，使 \(0<\abs{x-a}<\delta\) 时 \(\abs{g(x)-b}<\eta\)，直接代入即可。
\end{solution}

\section*{C组　结构、边界与项目}
\addcontentsline{toc}{section}{C组　结构、边界与项目}

\begin{solution}{14.19}
任意 \(\delta>0\) 可取 \(n>1/\delta\)，使 \(1/n\in D\) 且 \(0<1/n<\delta\)，故 \(0\) 是聚点。若极限为 \(L\)，偶数指标与奇数指标给出的值分别为 \(1,-1\)。取
\[
 \varepsilon_0=\frac12.
\]
若 \(L\ge0\)，在任意邻域选足够大的奇数 \(n\)，则 \(\abs{-1-L}\ge1\)；若 \(L<0\)，选足够大的偶数 \(n\)，则 \(\abs{1-L}>1\)。所有候选 \(L\) 均被排除。
\end{solution}

\begin{solution}{14.20}
设 \(\lim_{x\to a}f(x)=L\)，并给定 \(L\) 的任意邻域半径 \(\varepsilon>0\)。按定义存在 \(\delta>0\)，使所有满足
\[
 x\in D,\qquad0<\abs{x-a}<\delta
\]
的点都有 \(f(x)\in(L-\varepsilon,L+\varepsilon)\)。任何趋于 \(a\) 且最终不等于 \(a\) 的定义域点族，按“趋于”的含义最终进入上述 \(\delta\)-邻域，故对应函数值最终进入给定的 \(L\)-邻域。这里只是直接嵌套两个定义，没有使用逆向的序列刻画。
\end{solution}

\begin{solution}{14.21}
结论一般不成立。令 \(L=1\)，\(f(x)=1\)（\(x\in\Q\)）而 \(f(x)=-1\)（\(x\notin\Q\)），则 \(f^2(x)\equiv1\)，但 \(f\) 在任一点都没有极限。

一个充分条件是 \(f\) 在某删心邻域内保号。若最终 \(f(x)\ge0\)，则
\[
 \abs{f(x)-\abs L}
 =\frac{\abs{f^2(x)-L^2}}{f(x)+\abs L}
\]
在 \(L\ne0\) 时趋于零；若 \(L=0\)，由 \(\abs{f(x)}=\sqrt{f^2(x)}\to0\) 得结论。最终非正的情形类似，极限为 \(-\abs L\)。
\end{solution}

\begin{solution}{14.22}
恒等式
\[
 \max\{u,v\}=\frac{u+v+\abs{u-v}}2,\qquad
 \min\{u,v\}=\frac{u+v-\abs{u-v}}2
\]
把问题化为和、差、数乘以及绝对值的极限运算。若 \(f\to L\)、\(g\to M\)，便有
\[
 \max\{f,g\}\to\max\{L,M\},\qquad
 \min\{f,g\}\to\min\{L,M\}.
\]
\end{solution}

\begin{solution}{14.23}
可按以下四个模块整理。
\begin{enumerate}
 \item 因式分解：把误差写成 \(\abs{x-a}\) 与局部可控因子的乘积；平方、立方和多项式是典型对象。
 \item 有理化：用共轭式把根式差化为自变量差；中心值为零时需另行处理分母下界。
 \item 局部有界因子：由存在有限极限推出局部有界，再控制乘积误差；缺少有界性时，零因子乘无界因子可以产生多种行为。
 \item 分母保离零：非零极限给出统一正下界，从而允许倒数运算；分母极限为零时商法则失效。
\end{enumerate}
每个模块的证明都应先注明固定局部限制承担的作用，再按给定 \(\varepsilon\) 选择最终半径。
\end{solution}

\begin{solution}{14.24}
极限始终相对于定义域。以公式
\[
 f(x)=\frac{\abs x}{x}\qquad(x\ne0)
\]
为例：在 \(D_+=(0,\infty)\) 上令 \(x\to0\)，函数恒为 \(1\)，极限为 \(1\)；在 \(D_-=(-\infty,0)\) 上极限为 \(-1\)；在 \(D=\R\setminus\{0\}\) 上双侧极限不存在。公式相同，允许逼近中心的方向不同，量词所覆盖的点集随之改变。
\end{solution}
